\section{Exact Resources, Typed Algebras, and Programs}
\label{sec:cost-algebra-program}

This section describes the semantic spine used by finite cryptographic
computations.  Resource annotations are first-class results of execution, not
metadata reconstructed from an algorithm after the fact.  A typed primitive
handler is the only source of primitive costs, and a program obtains both its
ordinary distribution and its resource certificates from that same handler.
Figure~\ref{fig:exact-program-spine} summarizes the resulting separation.

\begin{figure}[t]
  \centering
  \resizebox{.97\linewidth}{!}{%
  \begin{tikzpicture}[node distance=8mm and 10mm]
    \node[rootbox] (model) {\decl{CostModel}\\ordered sequential resource};
    \node[rootbox, right=18mm of model] (signature)
      {\decl{Signature}\\result-indexed operations};

    \node[archbox, below=of model] (writer)
      {\decl{RandCosted}\(M\,X\)\\joint value--cost distribution};
    \node[archbox, below=of signature] (handler)
      {\decl{CostedAlgebra.exec}\\sole exact primitive handler};
    \node[archbox, below=of handler] (code)
      {\decl{Program.Code}\\\textsf{pure}, \textsf{bind}, \textsf{call}, \textsf{branch}};
    \node[archbox, below=of code] (run)
      {\decl{Program.Code.runCosted}\\exact program distribution};
    \node[rootbox, below=of run] (execution)
      {\decl{Program.Code.Execution}\\structural path and exact cost};
    \node[rootbox, below=of writer] (laws)
      {\decl{AlgebraLaws}\\cost-erased specification};
    \node[rootbox, below=of laws] (value)
      {\decl{Program.Code.valueDist}\\ordinary probability semantics};

    \node[certbox, right=18mm of handler] (opbound)
      {\decl{OperationBounds}\\per-operation certificate};
    \node[certbox, below=of opbound] (codebound)
      {\decl{Program.Code.Bound}\\structural certificate};
    \node[certbox, below=of codebound] (programbound)
      {\decl{BoundedProgram}\\one program, input-indexed budget};

    \draw[archarrow] (model) -- (writer);
    \draw[archarrow] (signature) -- (handler);
    \draw[archarrow] (writer) -- (handler);
    \draw[archarrow] (handler) -- (code);
    \draw[archarrow] (code) -- (run);
    \draw[<->,semithick,draw=cryptoteal] (run) --
      node[right,font=\scriptsize,align=left] {support\\equivalence} (execution);
    \draw[erasearrow] (run) -- node[above,sloped,font=\scriptsize]
      {erase cost} (value);
    \draw[erasearrow] (handler) -- (laws);
    \draw[certarrow] (handler) -- (opbound);
    \draw[certarrow] (opbound) -- (codebound);
    \draw[certarrow] (codebound) -- (programbound);
    \draw[certarrow] (programbound) -- (run);
  \end{tikzpicture}}
  \caption{The exact semantic spine and its proof-only side layers.  Blue
    arrows construct execution; dashed arrows forget costs; gold arrows build
    upper-bound evidence.  \decl{AlgebraLaws} and \decl{OperationBounds} do not
    define a second interpreter.  The two-headed edge is the proved
    correspondence between structural execution and the support of
    \decl{runCosted}; it does not identify the execution relation with
    \decl{valueDist}.}
  \label{fig:exact-program-spine}
\end{figure}

\subsection{Ordered sequential resource semantics}

The structure \decl{CostModel} packages a type \(C\), an additive monoid
\((C,0,+)\), a partial order \(\leq\), and monotonicity of addition in both
arguments.  Exact sequential execution uses \(+\) in evaluation order.  No
commutativity requirement is imposed, so a model may retain order-sensitive
information rather than merely count scalar steps.  The order and its two
monotonicity laws enter only when exact paths are compared with budgets.

The optional capability \decl{WorstCaseCostModel} adds an operation
\(\mathsf{sup}:C\times C\to C\) and proves its least-upper-bound laws against
the very same order stored by the base model.  It deliberately does not carry
a second order instance.  The concrete definitions \decl{CostModel.nat} and
\decl{WorstCaseCostModel.nat} instantiate sequential composition with natural
addition and worst-case choice with \(\max\), respectively; they are examples,
not assumptions of the generic semantics.

For a value type \(X\), \decl{Costed}\(M\,X\) is a pair
\(\langle x,c\rangle\) with \(x:X\) and \(c:C\).  Its writer bind is

\[
  \langle x,c_1\rangle \mathbin{\gg=} f
  = \langle y,c_1+c_2\rangle
  \quad\text{when}\quad f(x)=\langle y,c_2\rangle .
\]

The order of \(c_1+c_2\) is therefore observable for noncommutative models.
Theorems \decl{Costed.pure_bind}, \decl{Costed.bind_pure}, and
\decl{Costed.bind_assoc} establish the monad laws from the additive-monoid
laws; \decl{Costed.bind_cost} makes the single accumulation point explicit.

Randomness is added without marginalizing cost:

\[
  \mathsf{RandCosted}_M\,X \;:=\; \PMF(\mathsf{Costed}_M\,X).
\]

Thus a sampled value and the resource consumed on the path that produced it
remain correlated.  \decl{RandCosted.bind} combines PMF bind with writer bind,
and its lawful-monad proof includes associativity in the model's execution
order.  \mbox{\decl{RandCosted.valueDist}} is only the image under the value
projection.  In particular, the erasure theorem for bind is

\[
  \vDist(d \mathbin{\gg=} k)
    = \vDist(d) \mathbin{\gg=} (x \mapsto \vDist(k(x))),
\]

Within namespace \decl{RandCosted}, this is theorem \decl{valueDist_bind}.
The construction \decl{sampleWithCost} attaches a possibly value-dependent
cost to each sample; its value-erasure theorem proves that this annotation
leaves the sampled distribution unchanged.

The sole generic path-bound predicate is

\[
  \decl{RandCosted.CostBound}(d,b)
  \;\equiv\;
  \forall r\in\operatorname{supp}(d),\; r.\cost\leq b.
  \tag{1}
\]

Its \decl{pure}, \decl{map}, \decl{bind}, and \decl{weaken} theorems are the
basic proof algebra reused by later certificate layers.  In particular,
\decl{RandCosted.CostBound.bind} concludes a budget \(b_1+b_2\), in that
order, from support-wise bounds for the first computation and every
continuation.

\paragraph{Projection is observational, not semantic.}
The structure \decl{NatMeasure}\(M\) is a monotone additive homomorphism
\(\mu:C\to\Nat\).  It marks the boundary at which a rich exact resource is
observed as conventional runtime.  Theorems \decl{NatMeasure.map_add} and
\decl{NatMeasure.map_nsmul} cover sequential composition and finite
repetition.  The map \decl{RandCosted.mapCost} acts on every exact path and is
a writer-monad morphism; the value-erasure theorem
\mbox{\decl{RandCosted.valueDist_mapCost}} proves

\[
  \vDist(\decl{mapCost}\;\mu\;d)=\vDist(d).
\]

Monotonicity transports Eq.~(1) to the projected natural-number budget via
\decl{RandCosted.CostBound.mapCost}.  This projection therefore supports
complexity proofs while preserving both the underlying value distribution and
the unprojected exact run as the authoritative object.

\subsection{A typed algebra with separate laws and bounds}

An operation family often contains primitives with different result types:
sampling may return a scalar, scalar action returns a carrier element, and
addition returns another carrier element.  \decl{Signature} represents this
heterogeneity directly with a result-indexed family

\[
  \decl{Signature.Op} : (R:\mathsf{Type}) \to \mathsf{Type}.
\]

Consequently an inhabitant of \(\decl{S.Op}\,R\) carries, in its type, the fact
that handling it returns \(R\).  \decl{Signature.sum} forms the disjoint union
of two capabilities pointwise in the result index.  Its companion
\decl{CostedAlgebra.sum} dispatches to the corresponding exact handler, so
algorithms can request only the primitives they use rather than populate a
monolithic record with dummy fields.

The handler itself is intentionally small:

\[
  \decl{CostedAlgebra.exec} : \decl{S.Op}\,R \to
    \mathsf{RandCosted}_M\,R.
\]

It is the single source of both operation outcomes and their exact path costs.
Two structures describe independent facts about that handler.  First,
\decl{AlgebraLaws}\(A\) stores a cost-erased PMF specification and proves that
erasing \decl{A.exec} yields it.  Second, \decl{OperationBounds}\(A\) chooses a
budget for each operation and proves Eq.~(1) for every supported handler
result.  Neither structure executes an operation or stores a competing exact
cost.  Their respective \decl{sum} constructors preserve this separation
under signature composition.

\begin{table}[t]
  \centering
  \caption{Algebra and program objects, with the information each is allowed
    to own.}
  \label{tab:algebra-program-ownership}
  \small
  \begin{tabular}{@{}>{\raggedright\arraybackslash}p{.27\linewidth}
      >{\raggedright\arraybackslash}p{.31\linewidth}
      >{\raggedright\arraybackslash}p{.34\linewidth}@{}}
    \toprule
    Object & Owns & Does not own \\
    \midrule
    \decl{Signature} & typed operation family & PMFs, costs, or bounds \\
    \decl{CostedAlgebra} & exact joint operation handler & mathematical law or second cost table \\
    \decl{AlgebraLaws} & erased PMF specification and equality & execution semantics \\
    \decl{OperationBounds} & independently chosen operation budgets & exact operation costs \\
    \decl{Program} & one input-indexed typed body & a budget or duplicate bounded body \\
    \decl{BoundedProgram} & that program and a structural certificate & a second algorithm or interpreter \\
    \bottomrule
  \end{tabular}
\end{table}

The library supplies typed operation families for addition, negation,
subtraction, scalar action, multiplication, and sampling.  For example, within
namespace \decl{SampleOperation}, \decl{algebra} accepts an arbitrary exact
joint value--cost distribution.  Its \decl{laws} constructor separately
identifies the erased PMF, while \decl{bounds} separately certifies a path budget.
Deterministic arithmetic handlers follow the same pattern: their exact cost
functions occur only in the corresponding \decl{algebra} constructor, their
\decl{laws} constructors erase to the relevant Mathlib operation, and their
\decl{bounds} constructors compare exact costs with caller-selected budgets.

\begin{designprinciple}
Changing a mathematical specification cannot change execution, and changing
an upper-bound certificate cannot change either values or exact costs.  To
change an exact primitive cost, one must change the handler whose execution is
already used by every program consumer.
\end{designprinciple}

\subsection{One program, one interpreter, and support-complete executions}

For a fixed algebra \(A\), \decl{Program.Code}\(\,A\,R\) has four constructors:
\decl{pure} returns a value at zero cost; \decl{bind} sequences higher-order
code; \decl{call} invokes a typed primitive; and \decl{branch} selects between
two code values of the same result type using a Boolean condition.  External
input is not baked into this syntax.  Instead,

\[
  \mathsf{Program}\;A\;\mathit{Input}\;\mathit{Output}
  \quad\text{stores}\quad
  \decl{body}:\mathit{Input}\to\decl{Program.Code}\,A\,\mathit{Output}.
\]

This boundary permits the exact resource budget to depend on the supplied
input while keeping the program code structurally inductive.  Since
\decl{bind} is higher order, a primitive result may determine all subsequent
operations and the Boolean supplied to \decl{branch}.

The recursive definition \decl{Program.Code.runCosted} is the only evaluator.
It maps \decl{pure} to zero-cost writer pure, \decl{bind} to randomized writer
bind, \decl{call} to \decl{A.exec}, and \decl{branch} to the selected recursive
run.  \decl{Program.Code.valueDist} is definitionally the value projection of
this result.  Its compositional lemmas cover all four constructors, and
\decl{Program.Code.valueDist_call_eq} connects a call to its
\decl{AlgebraLaws} specification without adding an alternative evaluator.

The inductive relation \decl{Program.Code.Execution}\(
\mathit{code},\mathit{value},\mathit{cost}\) records
the actual structural path.  Its bind constructor adds the first and
continuation costs in order; its call constructor requires membership in the
support of \decl{A.exec}; and its two branch constructors record only the
selected branch.  Within namespace \decl{Program.Code}, the two directions are
proved separately and combined by
\decl{execution_iff_mem_support_runCosted}:

\[
  \decl{Execution}\;p\;x\;c
  \quad\Longleftrightarrow\quad
  \langle x,c\rangle\in\operatorname{supp}(\runC(p)).
  \tag{2}
\]

Equation~(2) is stronger than a one-way soundness statement: every semantic
support point has a structural witness, and every structural execution occurs
in the exact interpreter's support.  It also prevents a proof-only execution
relation from silently describing paths absent from the probabilistic
semantics.

\subsection{Structural and input-dependent bounds}

At code level, \decl{Program.Code.CostBound}\(
\mathit{code},\mathit{budget}\) is only an
abbreviation for Eq.~(1) applied to \decl{runCosted}\((\mathit{code})\).  The proposition
\decl{Program.Code.Bound}\(
\mathit{bounds},\mathit{code},\mathit{budget}\) is a structural certificate
indexed by the existing code; the code is not stored as another field.  Its
constructors mirror the semantic evaluator:

\begin{itemize}
  \item \decl{Bound.pure} certifies budget \(0\), and \decl{Bound.call} consumes
    the independently supplied \decl{OperationBounds};
  \item \decl{Bound.bind} combines budgets as \(b_1+b_2\) in execution order,
    requiring one common continuation budget valid for every intermediate
    value;
  \item \decl{Bound.branchOfBounds} accepts any caller-proved common upper
    bound for the two branch budgets; and
  \item \decl{Bound.branchSup} derives that common budget automatically only
    when the model provides \decl{WorstCaseCostModel}.
\end{itemize}

The distinction between the last two constructors is material.  Exact
interpretation needs no least-upper-bound operation.  Models without such an
operation remain usable, provided a caller supplies both inequalities to an
explicit common budget.  \decl{Bound.weaken} then permits a certified budget
to be widened using the model's partial order.

At the outer boundary,

\[
  \decl{Program.CostBound}\;p\;B
  \;\equiv\;
  \forall i,\;
    \decl{Program.Code.CostBound}\;(p.\decl{body}\;i)\;(B(i)),
\]

so budgets may depend on input.  \decl{BoundedProgram}\(
\mathit{bounds},\mathit{budget}\) stores
one \decl{Program} plus, for each input, a structural \decl{Code.Bound} for
that same body.  The theorem \decl{BoundedProgram.costBound} forgets the
structure to obtain the semantic path bound.  Its companion support corollary
applies that bound directly to an exact supported result.  Thus bounded and
unbounded views do not duplicate an algorithm: the former is the latter paired
with evidence about the exact interpreter already used for probability
semantics.

\begin{scopeboundary}
The program syntax charges only calls handled by the algebra's \decl{exec}
field.
Lean-level construction of higher-order continuations and evaluation of the
Boolean already supplied to \decl{branch} are not assigned an implicit host
runtime.  Adequacy of a chosen primitive cost model for a concrete execution
platform is therefore an external modeling obligation, not a theorem of this
semantic layer.
\end{scopeboundary}
